\section{A birational constancy criterion}
\label{cc:section}

Section~\ref{sec:rankone} made the entire highest line flat, with
finite character, along the relevant Iitaka fibers. We now turn this
Hodge-theoretic information into birational constancy. The argument is
first carried out over a curve, where a single lifted vector field can
be analyzed, and then spread to arbitrary algebraic parameter loci.

\begin{theorem}[Birational constancy over a curve]
\label{cc:criterion}
Let $\pi^\circ:P^\circ\to C^\circ$ be a smooth projective morphism with
connected fibers, where $C^\circ$ is a smooth connected complex algebraic
curve.  Suppose that its very general fiber $D$ satisfies
\[
                 \kappa(D)=0,\qquad h^0(D,K_D)=1.
\]
Put $n=\dim D$ and assume that the entire top Hodge bundle
$F^n(R^n\pi^\circ_*\C\otimes\cO_{C^\circ})$ is a flat line with
finite monodromy character.  Then there are a finite
extension $L/\C(C^\circ)$ and a smooth connected projective variety
$D_0/\C$ such that the generic fiber after extension to $L$ is birational
to $(D_0)_L$.

If a finite group $\Gamma$ acts on $P^\circ$ over $C^\circ$, one can choose
$D_0$ with a $\Gamma$-action and choose this birational identification
$\Gamma$-equivariantly.  The same conclusion holds for a finite
fiberwise birational action, after replacing the family by an
equivariant smooth projective model.
\end{theorem}

The proof lifts a base direction to a meromorphic vector field on the
total space.  We first control this field along the zero divisor of a
global top form.  A singular-metric estimate then shows that every pole
is a fixed divisor of a small ample perturbation of the canonical
system.  Contracting those divisors makes the vector field regular;
its flow identifies nearby fibers.  Finite changes of the curve and
birational changes of the total space over a dense open do not change
the asserted conclusion.  A zero-dimensional connected smooth fiber is
a point, so throughout the construction we assume $n>0$.

In the log-smooth K\"ahler setting of Cao--Guenancia--P\u aun,
with klt boundary and numerically trivial fiberwise log-canonical
class, flatness of the pluricanonical direct image in its
Narasimhan--Simha metric implies analytic local triviality of the
pairs over the smooth locus
\cite[\S1.1, Theorem~1.2]{CaoGuenanciaPaun23}.
Here the flat line is the entire top Hodge line, with finite
character, and the fibers satisfy $\kappa(D)=0$ and $h^0(D,K_D)=1$.
The top form may have zeros, allowing poles in the lifted base
vector field; contracting its possible poles yields the birational
conclusion.

\subsection{A global form and its polarization}

\begin{lemma}
\label{cc:preparation}
After a finite change of the curve there are a smooth projective curve
$C$, a smooth projective variety $P$ of dimension $n+1$, and a projective
morphism $\pi:P\to C$ whose restriction over a nonempty open is
birational to the changed family, with the following properties.
There is a closed holomorphic $n$-form $\Theta$ on $P$ whose restriction
to each fiber over that open spans its ordinary top-form space.  Write
\[
 s\in H^0(P,K_{P/C}),\qquad \mathcal E=\operatorname{div}(s)
\]
for its relative canonical section and zero divisor.  The horizontal
part of $\mathcal E$ is relatively simple normal crossing over a smaller
open of $C$: the morphism and every horizontal zero-divisor stratum
are smooth there.  We call fibers over this open \emph{good fibers}.

One can choose an ample line bundle $H$ with a positive smooth metric,
whose curvature form is denoted by $h$, so that, for a positive area
class $\ell$ pulled back from $C$, there is a holomorphic $n$-form
$\alpha$ satisfying
\begin{equation}
           [h\wedge\Theta]=[\ell\wedge\alpha]
                  \quad\text{in }H^{n+1,1}(P).
\label{cc:pairing}
\end{equation}
In the equivariant case $P$ and $H$ can be chosen equivariant, $h$ is
invariant, and $\Theta$ transforms by the character of the fiber top
line.  Moreover, $P$ dominates a smooth projective model $P_0\to C$
with reduced fibers, and $\Theta$ is the pullback of a holomorphic form
on $P_0$.
\end{lemma}

Identity~\eqref{cc:pairing} has two uses: it removes the obstruction
to lifting the base vector field constructed below to $H$, and
later converts a weighted
integral on $P$ into an integral over good fibers. We obtain it by
choosing the global lift of the flat top form; that choice will leave
the relative form $s$ and its zero divisor unchanged.

\begin{proof}
Kill the finite character by a finite cover of $C^\circ$ and complete
the resulting curve.  Semistable reduction over a curve, followed by the
usual toroidal resolution, gives, after a further finite change, a smooth
projective model $P_0\to C$ with reduced fibers; see
\cite[Chapter II, semistable reduction over a curve]{KKMSD}.
Take a smooth projective common dominating model of this family and the
original changed family.  When $\Gamma$ is present, include the finitely
many translates in the graph construction and use an equivariant
resolution.  This also regularizes a finite birational action: take the
closure of the graph of all its maps in the product of their projective
models, on which the group permutes the factors, and resolve
equivariantly.  Over a dense smooth open, birational pullback preserves
ordinary top forms and is compatible with the cohomology local systems.
The chosen top class is therefore still invariant.

The flat generator already gives a relative top form on this smooth
open.  Choose a log resolution of its geometric generic zero divisor.
That resolution and its finitely many geometric components are defined
after a finite extension of the curve's function field.  Make this
change now and spread the resolution.  Repeat semistable preparation
at the boundary and take a smooth common dominating model, preserving
the chosen resolution over a dense open.  Thus the geometric divisor
components used below are defined before the final $P$, $H$, and
$\Theta$ are fixed.

The global invariant cycle theorem gives a surjection from the
cohomology of the smooth compactification onto the invariant part of
the fiber cohomology.  The restriction map is a morphism of Hodge
structures, so its strictness gives a preimage of the invariant
$(n,0)$-class in $H^{n,0}(P)$; these are precisely holomorphic $n$-forms.
Here we use \cite[Theorem~4.1.1(ii) and Corollary~4.1.2]{Deligne71}, with
smooth proper family over the good curve and smooth projective total
compactification.  Holomorphic forms on a smooth projective variety are
closed.

To describe the relative section, choose a local coordinate $t$ on $C$.
The absolute form $dt\wedge\Theta$, divided by the base frame $dt$,
defines $s$ in $K_{P/C}$.  These expressions agree under a change of
coordinate on the curve.  The section is holomorphic globally and is
nonzero on every good fiber.  Resolve its horizontal zero divisor and
shrink the good open so that the resulting horizontal components and
their intersections are smooth over the curve.  This gives the stated
relative normal crossings.  Any two lifts of the fixed flat top class
have the same restriction to every good fiber, and therefore give the
same relative section.  Thus the following change of lift does not alter
$s$ or this preparation of its zero divisor.

Put $V=H^0(P,\Omega_P^n)$ and let $N$ be the kernel of restriction to a
good fiber.  This kernel is independent of the chosen good fiber:
restriction of a global closed form gives a flat cohomology section, and
a holomorphic top form represents the zero cohomology class only when
it is zero.  Choose an ample $H$ and a positive metric, and set
\[
 B_h(\beta,\gamma)=\mathrm{i}^{n^2}
             \int_P h\wedge\beta\wedge\overline\gamma.
\]
This is a positive definite Hermitian pairing on $V$.  Indeed its
integrand is a positive multiple of the pointwise norm squared when
$\beta=\gamma$, and a nonzero holomorphic form is nonzero on an open
set.  For a positive smooth area form $\ell$ from the curve, the pairing
$B_\ell$ is semidefinite.  Fiber integration identifies its kernel with
$N$.  Consequently the image of the map
\[
       V\longrightarrow\overline V^{\,*},\qquad
       \beta\longmapsto B_\ell(\beta,-)
\]
is the annihilator of $N$.

Choose the lift $\Theta$ orthogonal to $N$ for $B_h$.  The functional
$B_h(\Theta,-)$ annihilates $N$, so it equals $B_\ell(\alpha,-)$ for
some $\alpha\in V$.  The duality pairing of $H^{n+1,1}(P)$ with
$H^{0,n}(P)$ is nondegenerate.  Applying it to the difference proves
\eqref{cc:pairing}.

In the equivariant case a tensor product of the translates of an ample
bundle admits a natural $\Gamma$-linearization; its metric can be made
invariant while keeping positive curvature.  Both $N$ and the pairing
are then invariant.  The unique orthogonal lift of the chosen top class
transforms by its top-line character.  The same character can be imposed
on $\alpha$ by projection to its character space.

Finally, a holomorphic form on a smooth projective birational model is
the pullback of a holomorphic form on every smooth projective model it
dominates.  Applying this to $P\to P_0$ gives the asserted descent of
$\Theta$.  A later positive tensor power of $H$ multiplies $h$ by a
positive integer and preserves all these conclusions, with $\alpha$
scaled accordingly.
\end{proof}

On the complement of the zeros and the critical locus, the kernel of
$\Theta$ is a holomorphic line transverse to the fibers.  On a good base
coordinate disk it has a unique generator $v$ satisfying
\begin{equation}
        d\pi(v)=\partial_t,\qquad \iota_v\Theta=0.
\label{cc:kernel}
\end{equation}
Here and below we use $\mathcal E$ for a divisor on the total space and
$E=\mathcal E|_D$ for its divisor on a good fiber.

\begin{lemma}[A meromorphic differential operator]
\label{cc:operator}
The vector field $v$ extends meromorphically to a neighborhood of each
general fiber, with poles bounded by $\mathcal E$.  After shrinking the
base disk, it is the symbol of a meromorphic first-order differential
operator $\nabla$ on $H$ with the same pole bound.
\end{lemma}

\begin{proof}
In a local relative canonical frame, write $\omega$ for the relative
$n$-form represented by $s$.  The normalization in
\eqref{cc:kernel} gives
\begin{equation}
                 \Theta=\iota_v(dt\wedge\omega).
\label{cc:contraction}
\end{equation}
Thus $sv$ is a holomorphic section of $T_P\otimes K_{P/C}$: contraction
with an absolute top form identifies this bundle, after choosing $dt$,
with $\Omega_P^n$.  Dividing by $s$ proves the pole bound.

Use the absolute scalar-symbol sequence
\begin{equation}
 0\longrightarrow\cO_P\longrightarrow\operatorname{At}(H)
                   \longrightarrow T_P\longrightarrow0.
\label{cc:atiyah}
\end{equation}
Its middle sheaf consists of first-order differential operators on $H$
with scalar symbol.  Its extension class is $c_1(H)$, up to the fixed
normalizing constant; this is the connection obstruction of
\cite{Atiyah57}.  For the specific assertion needed here, take frames
$e_i=g_{ij}e_j$.  Local lifts of the same vector symbol differ by its
contraction with $d\log g_{ij}$.  This is the cocycle representing the
stated extension class.

Tensor \eqref{cc:atiyah} by $K_{P/C}$.  The obstruction to lifting $sv$
is its contraction with $c_1(H)$ in $H^1(K_{P/C})$ on the neighborhood
in question.  In Dolbeault notation and using
\eqref{cc:contraction}, its restriction to $D_t$ is, up to a nonzero
constant and sign, the class obtained from $h\wedge\Theta$ by dividing
its top holomorphic degree by $dt$ and restricting to $D_t$.
This operation respects Dolbeault coboundaries.  In fact the total space
has dimension $n+1$, so a primitive of type $(n+1,0)$ necessarily contains
$dt$ in submersion coordinates; division by $dt$ commutes with
$\bar\partial$ before restriction.  Equivalently this is the canonical
restriction $K_P|_{D_t}\simeq K_{D_t}\otimes T_t^*C$ with a base frame
chosen.

The right side of \eqref{cc:pairing} gives zero under this operation:
a smooth area representative of $\ell$ contains $d\bar t$.  Hence the
fiber obstruction is zero for every good $t$.  Shrink to a Stein disk on
which cohomology and base change hold for $K_{P/C}$.  Leray and the
vanishing of higher coherent cohomology on the disk identify the
obstruction with a section of $R^1\pi_*K_{P/C}$.  It is zero because its
fiber values are zero.  Therefore $sv$ lifts to a holomorphic section of
$\operatorname{At}(H)\otimes K_{P/C}$ on the whole neighborhood.
Dividing this lift by $s$ produces $\nabla$, with the claimed symbol and
pole bound.  Replacing $H$ by a tensor power uses the induced operator
and leaves its symbol and its pole bound unchanged.
\end{proof}

\subsection{Vanishing on the zero divisor}

The lifted field may have poles along $\mathcal E$.  We first prove
that the absolute form vanishes on every component of this divisor.
This will make its pairing with every positive curvature current of
$K_{P/C}$ vanish, and supply the integral estimate used below.

\begin{lemma}
\label{cc:zeros}
The pullback of $\Theta$ to a resolution of every irreducible component
of $\mathcal E$ is zero.
\end{lemma}

\begin{proof}
We first treat a horizontal component.  Choose a simultaneous very
general good fiber $D$, so that all its positive pluricanonical section
spaces have dimension one.  Write
\begin{equation}
 E=\sum_j e_jD_j,\qquad R_0=E+E_{\mathrm{red}},\qquad
 S=D_i,\qquad B'=\sum_{j\ne i}D_j.
\label{cc:zero-data}
\end{equation}
The coefficients $e_j$ are positive integers, $E\sim K_D$, and the
support is simple normal crossing.  The components extend relatively
over a small disk.  In the local computations we use the same letters
for these relative extensions and their local equations.

The contradiction will take place in the section spaces of $R_0$.
Since
\[
                   E\le R_0\le2E,\qquad E\sim K_D,
\]
for every positive integer $m$ we have
\[
       1\le h^0(D,mR_0)\le h^0(D,2mE)
                              =h^0(D,2mK_D)=1.
\]
Thus every section of $mR_0$ is a multiple of its divisor section
and vanishes on $S$. We will construct a section whose restriction
to $S$ is nonzero.

Suppose that the restriction of $\Theta$ to the horizontal component
through $S$ is nonzero.  Let $z$ and $G$ be local equations for $S$ and
$R_0$, and put
\[
                   W=Gv,\qquad u=(W(z)/z)|_S.
\]
The pole bound in Lemma~\ref{cc:operator} shows that $W$ is holomorphic
and vanishes on the reduced zero divisor.  It is consequently
logarithmic along that divisor, so that both $W(z)/z$ and $W(G)/G$ are
holomorphic.  Under changes of the equations, $u$ transforms as a
section of $R_0|_S$.  Formula \eqref{cc:contraction} identifies its
nonvanishing at the general point of $S$ with the supposed nonvanishing
of $\Theta$ on the relative component.  Thus
\begin{equation}
                0\ne u\in H^0(S,\cO_S(R_0)).
\label{cc:u}
\end{equation}
Moreover $u$ is divisible by the full reduced divisor $B'|_S$.  To see
this at a crossing, the coefficients of $sv$ are holomorphic, while
$G/s$ has one factor for each component of $E_{\mathrm{red}}$.
Division by $z$ removes the factor belonging to $S$ and leaves all
other reduced factors.

Write the operator in a regular frame of $H$ as $\nabla=v+c$.  Its pole
bound also gives $Gc|_S=0$.  Replace $H$ by a sufficiently high tensor
power.  Relative generation supplies finitely many sections
$a_\lambda$ on a neighborhood of the whole fiber, vanishing on the
relative $S$, such that $a_\lambda/z$ generate $H(-S)$ along $S$.
For each $\lambda$ set
\[
 q_{k,\lambda}=\frac{\nabla^k a_\lambda}{k!},\qquad
 b_{k,\lambda}=\frac{G^k\nabla^ka_\lambda}{z}.
\]
The following recurrence proves in particular that the second
expression is holomorphic:
\begin{align}
 b_{k+1,\lambda}
  &=W(b_{k,\lambda})+
     \left(\frac{W(z)}z-k\frac{W(G)}G+Gc\right)b_{k,\lambda},
                                                       \label{cc:recurrence}\\
 b_{k,\lambda}|_S
  &=u^k(a_\lambda/z)|_S
           \prod_{j=0}^{k-1}\bigl(1-j(e_i+1)\bigr).
                                                       \label{cc:leading}
\end{align}
For the second equality, $W$ vanishes on $S$, $Gc|_S=0$, and
$(W(G)/G)|_S=(e_i+1)u$.  The other components of $G$ contribute zero
on $S$; this can be checked away from their intersections and then
extended holomorphically.  These facts prove both formulas by induction.
They also show that the restrictions of $q_{k,\lambda}$ are sections of
$\cO_D(kR_0+H-S)$. More explicitly, in a local frame of $H$,
\[
 q_{k,\lambda}=\frac{z}{G^k}\frac{b_{k,\lambda}}{k!}.
\]
The meromorphic expression on the left is therefore a holomorphic
section of the indicated twisted bundle, with local coefficient
$b_{k,\lambda}/k!$.

We need estimates as $k$ tends to infinity.  Cover a compact
neighborhood of $D$ by finitely many pairs of nested coordinate charts.
For a fixed $k$, divide each nesting margin into $k$ equal parts.  On
these charts the coefficients of $W$, $W(z)/z$, $W(G)/G$, and $Gc$ are
bounded.  At the $j$th step of \eqref{cc:recurrence}, Cauchy's estimate
for the derivative costs at most a constant times $k$, and the
zeroth-order coefficient costs at most a constant times $1+j\le k+1$.
It follows that the holomorphic coefficients $b_{k,\lambda}/k!$, in
smooth metrics for $kR_0+H-S$, are bounded by
\begin{equation}
                   \frac{(Ck)^k}{k!}\le C_1^k.
\label{cc:cauchy}
\end{equation}
This is a bound in the indicated twisted bundle, including across the
zero divisor.

Put $a=e_i+1\ge2$.  No factor in \eqref{cc:leading} is zero, and
\[
 \lim_{k\to\infty}
 \left|\frac{\prod_{j=0}^{k-1}(1-ja)}{k!}\right|^{1/k}=a>0;
\]
for example, the ratio of successive absolute values tends to $a$.
On $D$, take the weights
\[
 \frac1k\log\sum_\lambda
       \left|\frac{b_{k,\lambda}}{k!}\right|^2
       -\frac1k\varphi_{H-S},
\]
where $\varphi_{H-S}$ is a smooth weight on $H-S$ in the same
local frames. They are weights on $R_0$ with curvature loss tending
to zero and uniform upper bounds by \eqref{cc:cauchy}.  Their
regularized upper limit is a semipositive singular weight.  Formula
\eqref{cc:leading}, together with generation of $H(-S)|_S$, gives its
restriction a lower bound by $\log|u|^2$ up to a constant.  Finally take
the maximum with the divisor weight of $R_0$.  We obtain a semipositive
weight $\varphi$ on $R_0$ satisfying
\begin{equation}
      \varphi\ge\log|s_{R_0}|^2,\qquad
      \varphi|_S\ge\log|u|^2-C.
\label{cc:weight-bounds}
\end{equation}
All weights here are written in compatible local frames; the
inequalities are invariant statements about the corresponding metrics.

We now extend a nonzero power of $u$ using
\cite[Theorem~4.1, conditions~(19)--(23)]{DHP13}.  We record every
hypothesis in this application.  The ambient manifold is the smooth
projective $D$, and its smooth hypersurface is $S$.  For any integer
$m\ge2$ put
\begin{equation}
 \begin{split}
  F'&=mR_0-(K_D+S)\sim (m-1)R_0+B',\\
  \varphi_{F'}&=(m-1)\varphi+\log|s_{B'}|^2,\\
  A_0&=R_0-aS=\sum_{j\ne i}(e_j+1)D_j,\\
  \varphi_S&=1+\frac{\varphi-\log|s_{A_0}|^2}{a}.
 \end{split}
\label{cc:extension-weights}
\end{equation}
The first line uses $K_D\sim E$ and
$E_{\mathrm{red}}=S+B'$.  Thus $F'$ is an actual line bundle with the
indicated weight, after the indicated line-bundle identification.

Condition (19) of the extension theorem holds with its constant equal
to one, since \eqref{cc:weight-bounds} implies
\[
                  |s_S|^2e^{-\varphi_S}\le e^{-1}.
\]
For condition (20), the weights $1+\varphi/a$ and
$a^{-1}\log|s_{A_0}|^2$ are semipositive weights on rational line
bundles and have difference $\varphi_S$.  If integral bundles are used,
choose an integral $G_2$ sufficiently ample that $G_2-A_0/a$ has a
positive smooth rational-bundle metric, and add this metric to both
weights.  The difference is unchanged and both underlying bundles
become integral.

For condition (21), let $T_R=\ddc\varphi$.  Directly,
\begin{equation}
 \begin{split}
 \ddc\varphi_{F'}&=(m-1)T_R+[B']\ge0,\\
 \ddc\varphi_{F'}-\ddc\varphi_S
       &=\left(m-1-\frac1a\right)T_R+[B']+\frac1a[A_0]\ge0.
 \end{split}
\label{cc:extension-curvature}
\end{equation}
Here $a\ge2$ and $m\ge2$.  For condition (22), choose a positive
number
\[
  \varepsilon_0\le
       \min\left\{1,\min_{j\ne i}\frac{a}{e_j+1}\right\},
\]
omitting the inner minimum if there are no other components.  The
weight $\varphi$ is locally bounded above.  After normalizing local
defining functions to have norm at most one, each coefficient of
$\log|s_{B'}|^2$ is at least the corresponding coefficient in
$\varepsilon_0 a^{-1}\log|s_{A_0}|^2$.  Therefore
\[
       \varphi_{F'}\le
             \frac{\varepsilon_0}{a}\log|s_{A_0}|^2+C.
\]
A common smooth metric added to the two weights above changes this
inequality only by a bounded local term.  This proves precisely the
required domination of the subtracted weight's singularities.

The section to extend is
\[
       u^m\in H^0\bigl(S,\cO_S(mR_0)\bigr)
             =H^0\bigl(S,K_S+F'|_S\bigr).
\]
The lower bound for $\varphi|_S$ and $u\ne0$ also ensure that
the restricted weights are not identically minus infinity.
Condition (23) follows from the second bound in
\eqref{cc:weight-bounds}:
\begin{equation}
 |u|^{2m}e^{-\varphi_{F'}|_S}
                \le C\frac{|u|^2}{|s_{B'}|_S|^2}.
\label{cc:extension-integral}
\end{equation}
The quotient on the right is locally bounded because $u$ is divisible
by the full reduced divisor $B'|_S$.  There is in fact integrability
with a slightly larger weight exponent.  Writing $u=s_{B'}w$, the
density with exponent $1+\delta$ on $\varphi_{F'}$ is bounded by
\[
 C|s_{B'}|^{-2m\delta}|w|^{2-2(m-1)\delta}.
\]
It is integrable for $0<\delta<1/m$, by the normal crossings of
$B'|_S$ and the positive exponent on $|w|$.  This argument requires no
normal-crossing assertion about the other zeros of $w$.

The extension theorem now gives a section of $mR_0$ with restriction
$u^m\ne0$, contradicting the one-dimensionality of this section
space established above: its divisor section vanishes on $S$.
This proves the horizontal vanishing.

For vertical components, use the reduced-fiber model $P_0\to C$ in
Lemma~\ref{cc:preparation}, to which $\Theta$ descends.  At a general
point of a reduced fiber component the map has local equation $t=z$.
There the vanishing of the relative coefficient $s$ is exactly the
vanishing of the restriction of $\Theta$ to that component.  Every
vertical component on $P$ which is not exceptional over $P_0$ is tested
at the same generic point.  An exceptional divisor maps into a subset
of $P_0$ of dimension at most $n-1$, and the pullback of an $n$-form to
its resolution is zero.  These observations prove the remaining
vanishing assertions.  They also explain why the reduced-fiber model
is used: for $t=z^r$ with $r>1$, an additional differential factor
$z^{r-1}$ would not itself imply vanishing of the absolute form.
\end{proof}

\subsection{Curvature pairing and an integrability estimate}

We have proved that $\Theta$ vanishes on every component of
$\mathcal E=\operatorname{div}(s)$. Put
$\mu=i^{n^2}\Theta\wedge\overline\Theta$.
The resulting identity $[\mathcal E]\wedge\mu=0$ will control
singular metrics on $K_{P/C}$, and hence the poles of the normalized
kernel field $v$.

For a semipositive weight $T$ on $K_{P/C}$, the curvature pairing
with $\mu$ has the same integral as $[\mathcal E]\wedge\mu$ and
therefore vanishes. When $T\geq\log|s|^2$, applying this observation
to truncations shows that $\psi=T-\log|s|^2$ is constant along
almost every local kernel leaf. This permits the weighted use of
\eqref{cc:pairing} and gives an integrability exponent strictly
greater than one. We prove the estimate for arbitrary $T$ before
constructing the weight supplied by pluricanonical section spaces.

We use weights for metrics written as \(e^{-\varphi}\), and the convention
\[
 \ddc=\frac{i}{2\pi}\partial\bar\partial,
 \qquad \ddc\log|z|^2=[z=0].
\]
In particular, if \(\varphi_H\) is a local weight for the metric on
\(H\), then \(\ddc\varphi_H=h\).  Expressions such as
\(\log|s|^2\) below denote local weights in holomorphic frames; the
difference of two weights on the same bundle is a globally defined
function wherever both are finite.

\begin{lemma}\label{ca:integrability}
Let \(T\) be a semipositive singular weight on \(K_{P/C}\) such that
\(T\geq\log|s|^2\), and put
\(\psi=T-\log|s|^2\geq0\) off \(\mathcal E\).
There exists \(\delta>0\) such that
\begin{equation}\label{ca:global-integral}
 \int_P e^{(1+\delta)\psi}
       i^{n^2}h\wedge\Theta\wedge\overline\Theta<\infty.
\end{equation}
\end{lemma}

\begin{proof}
The form \(\mu\) is smooth, closed, and positive of type \((n,n)\).
Lemma~\ref{cc:zeros} implies that
\([\mathcal E]\wedge\mu=0\): the integral over each component can be
computed on a resolution of that component, where the pullback of
\(\Theta\) vanishes.  Since the curvature current of
\(\log|s|^2\) is \([\mathcal E]\), cohomology and closedness give
\[
 \int_P (\ddc T)\wedge\mu
    =\int_P[\mathcal E]\wedge\mu=0.
\]
Here and subsequently \(\ddc T\) denotes the globally defined curvature
current of the bundle weight.  The measure on the left is positive,
so it vanishes.  For every real \(b\), the weight
\[
 T_b=\max\{T,\log|s|^2+b\}
\]
is another semipositive weight on the same bundle.  The same argument
shows that \((\ddc T_b)\wedge\mu=0\).

Delete \(\mathcal E\), the zeros of \(\Theta\), and the bad-fibration
locus.  In a holomorphic flow box for the kernel of \(\Theta\), use
coordinates \((w_1,\ldots,w_n,\tau)\), where the leaves vary only in
\(\tau\).  Closedness gives
\[
 \Theta=a(w)\,dw_1\wedge\cdots\wedge dw_n,
 \qquad a(w)\ne0.
\]
The weight \(\log|s|^2\) is pluriharmonic here.  Thus wedging the
curvature of \(T\), or of \(T_b\), with \(\mu\) selects the
\(\tau\)-Laplacian of \(\psi\), or of \(\max\{\psi,b\}\), against
a nonvanishing transverse volume.  Slicing these zero positive
measures shows that both functions are harmonic on almost every
local leaf.  Take the intersection of these full-measure sets for
rational \(b\).  If a harmonic function is nonconstant, there is a
small leaf neighborhood on which its differential is nonzero and
whose image contains an open real interval.  Choose a rational
\(b\) in that interval.  Taking the maximum with \(b\) produces a
nonzero positive Laplacian on the level crossing in that
neighborhood.  Hence \(\psi\) is constant on almost every local
leaf.

In these coordinates $\overline\Theta$ contains all the transverse
differentials $d\bar w_j$. Thus only a derivative in the leaf
direction can contribute to
$\bar\partial(\chi(\psi)\overline\Theta)$, and leafwise constancy
gives, for every bounded continuous function \(\chi\),
\[
 \bar\partial\bigl(\chi(\psi)\overline\Theta\bigr)=0
\]
off the deleted analytic set, in the sense of distributions.  This
identity extends across that set.  Indeed, the coefficients of this
current are locally bounded.  Across a smooth analytic stratum of
complex codimension \(q\geq1\), a cutoff at distance \(\epsilon\)
has derivative \(O(\epsilon^{-1})\) on a tube of volume
\(O(\epsilon^{2q})\), so its contribution tends to zero.  Applying
this successively to the smooth strata, and choosing the tubes about
the remaining lower-dimensional strata sufficiently small, proves
the distributional extension.  This argument uses boundedness of
\(\chi\); it makes no differentiability assumption on \(\psi\).

We may now pair \eqref{cc:pairing} with
\(\chi(\psi)\overline\Theta\).  Choose the representative \(\ell\)
there to be the pullback of a nonnegative smooth area form supported
in relatively compact good base charts.  This does not change its
cohomology class.  On those charts, after shrinking within the good
locus, fiberwise rank one gives
\[
 \alpha|_{D_t}=a(t)s|_{D_t}
\]
with \(a(t)\) bounded on the support in question.  Dolbeault
integration by parts against the closed current just constructed
therefore gives
\begin{equation}\label{ca:paired-truncation}
 \int_P \chi(\psi)\,h\wedge\Theta\wedge\overline\Theta
 =\int_P\chi(\psi)\,\ell\wedge\alpha\wedge\overline\Theta.
\end{equation}
On a good fiber write \(\omega=s|_{D_t}\).  In local frames the
absolute value of the right-hand integrand for
\(\chi(x)=\min\{M,e^{(1+\delta)x}\}\) is bounded by a fixed multiple of
\[
 |\omega|^2 e^{(1+\delta)(T-\log|\omega|^2)}
   =e^{(1+\delta)T}|\omega|^{-2\delta}.
\]
Here \(|\omega|^2\) denotes the canonical volume density.
The weights \(T\) have local upper bounds.  The zero divisor of
\(\omega\) is relatively simple normal crossing on this compact
set, so the last expression is integrable if \(\delta>0\) is small
enough that \(\delta e_j<1\) for its finitely many multiplicities.
The bound is independent of \(M\).  Multiply
\eqref{ca:paired-truncation} by \(i^{n^2}\), take absolute values on
the right, and let \(M\to\infty\).  Monotone convergence proves
\eqref{ca:global-integral}.
\end{proof}

\subsection{Fixed orders and Bergman weights}

Calculate the following orders first on the geometric generic fiber.
The section-space and vanishing-filtration comparisons below will
give the same values on every simultaneous very general fiber.
For a prime divisor \(S\) on such a fiber \(D\) and positive rational
\(\epsilon\), define its actual asymptotic fixed order by
\[
 a_S(\epsilon)=
 \inf_{\substack{m>0\\m\epsilon\in\Z}}
 \frac{1}{m}\min_{0\ne U\in
 H^0(D,mK_D+m\epsilon H|_D)}\ord_S(U),
\]
where degrees with zero section space are omitted.  These spaces
are nonzero in sufficiently divisible degree, since \(K_D\sim E\)
and \(\epsilon H|_D\) is ample.  Multiplication by an ample
section not vanishing generically on \(S\) shows that
\(a_S(\epsilon)\) is nondecreasing as \(\epsilon\) decreases.
In making this comparison, first pass to common divisible degrees
and raise any chosen section to a power.  Thus the limit
\[
 \sigma_S(K_D)=\lim_{\epsilon\downarrow0}a_S(\epsilon)
\]
exists. We need only this definition and its elementary
monotonicity.

Our objective is to prove that $\sigma_S(K_D)>0$ whenever
$S$ comes from a horizontal pole divisor of $v$. A zero value
would give section degrees $k$ and positive ample twists $r_k$
for which both $r_k/k$ and the normalized minimal vanishing
order $j_k/k$ tend to zero. The next two lemmas associate to
such degrees a semipositive weight with zero generic divisorial
coefficient along $S$. The pole-density estimate will then show
that this weight cannot satisfy Lemma~\ref{ca:integrability}
at a pole of $v$.

The weight must be constructed on $P$ before we select the fiber
on which its global integral is finite. We therefore use relative
section spaces on the fixed model of Lemma~\ref{cc:preparation}.
The weight will work on every fiber satisfying the corresponding
base-change conditions. Only afterwards, in
Proposition~\ref{ca:poles}, will we make the final fiber choice.

For any good fiber \(D\) under consideration, write
\[
 \omega=s|_D,\qquad E=\mathcal E|_D=\sum_j e_jD_j,
 \qquad \varphi_0=\log|\omega|^2.
\]
The geometric divisor components are already defined on the fixed
model of Lemma~\ref{cc:preparation}.  Choose the integer degrees and
relative section spaces on this model; no individual sequence of
generic sections needs to be spread.  For each degree, cohomology and
base change also apply to the section
spaces with prescribed vanishing along these components.  Excluding
the resulting countable union of proper closed subsets fixes all
section dimensions and minimal vanishing orders simultaneously.

\begin{lemma}\label{ca:bergman}
Choose integers \(k\to\infty\) and \(r_k>0\) with
\(r_k/k\to0\), such that
\(H^0(D,kK_D+r_kH|_D)\ne0\) on the generic fiber.  There is a
semipositive weight \(T\) on \(K_{P/C}\), depending only on these
degrees and the fixed family and metrics, with \(T\geq\log|s|^2\).
On every good fiber \(D\) satisfying the simultaneous base-change
conditions for these unfiltered section spaces, for every sequence
\(U_k\in H^0(D,kK_D+r_kH|_D)\) normalized by
\begin{equation}\label{ca:normalization}
 \int_D |U_k|^{2/k}e^{-r_k\varphi_H/k}=1,
\end{equation}
the weights and functions
\begin{equation}\label{ca:section-weights}
 \varphi_k=\frac{\log|U_k|^2-r_k\varphi_H}{k},
 \qquad x_k=\varphi_k-\varphi_0,
 \qquad \psi=T|_D-\varphi_0
\end{equation}
satisfy uniform local upper bounds for \(\varphi_k\) and
\(\limsup_k x_k\leq\psi\) off \(E\).  The integral in
\eqref{ca:normalization} is understood in canonical coordinate
densities, and is independent of those coordinates and bundle
frames.
\end{lemma}

\begin{proof}
Apply the relative \(k\)-Bergman construction to
\(\pi:P\to C\) and the twisting bundle \(r_kH\) with its smooth
positive metric.  The source and base are smooth, the morphism is
projective and surjective with connected fibers, and a nonzero
fiber section exists generically.  Generic base change supplies a
nonempty open on which all such sections extend locally.  The
smooth twisting metric makes every fiber section \(2/k\)-integrable.
Thus Theorem~4.2.7 of \cite{PaunTakayama18}, in the cited arXiv version, supplies
a semipositive weight \(\beta_k\) on \(kK_{P/C}+r_kH\).  Over the
smooth extendibility locus, its exponential is the extremum of
\(|U|^2\) among sections normalized as in
\eqref{ca:normalization}.  No extremal interpretation on exceptional
fibers is needed.

Set
\[
 \tau_k=\frac{\beta_k-r_k\varphi_H}{k}.
\]
These are weights on \(K_{P/C}\), and
\(\ddc\tau_k\geq-(r_k/k)h\).  On a fixed smaller submersion chart,
the mean inequality for the subharmonic function given by the
absolute value of a section coefficient to the power \(2/k\),
together with its normalized integral and the bounded factor
\(e^{-r_k\varphi_H/k}\), gives an upper bound independent of \(k\).
Taking the extremum gives that bound for \(\tau_k\).  The bound
persists across degree-dependent exceptional base values by
plurisubharmonic extension.  Equivalently, the ambient estimate in
Remark~4.3.1 of \cite{PaunTakayama18} gives precisely this uniformity when the
divided twisting weights are bounded above.

For completeness, local boundedness on one fixed nonempty open also
globalizes as follows.  Subtract a fixed smooth reference weight on
\(K_{P/C}\).  The resulting globally defined potentials have a
uniform lower bound for their complex Hessians.  If their suprema
were unbounded, subtracting those suprema and using compactness of
sup-normalized quasi-plurisubharmonic potentials would give a
subsequence converging in \(L^1\) to a finite quasi-plurisubharmonic
function.  On the fixed open, however, the normalized functions
would tend uniformly to minus infinity.  This is impossible.
Hence the weights \(\tau_k\) are uniformly bounded above in local
smooth frames on \(P\).

These constructions use the relative section spaces, without selecting
a final fiber or an individual section.  Take their regularized upper
limit on $P$ and then its maximum with the divisor weight:
\begin{equation}\label{ca:limiting-weight}
 T=\max\left\{\left(\limsup_{k\to\infty}\tau_k\right)^*,
                         \log|s|^2\right\}.
\end{equation}
The curvature losses tend to zero, so this is semipositive; if the
first term is identically minus infinity, the second term still
defines the weight.  Its restriction to a good fiber is well
defined because it dominates \(\varphi_0\).  Every normalized
section is bounded above by the corresponding Bergman extremum.
Consequently \(\varphi_k\leq\tau_k|_D\), which proves all the
assertions.
\end{proof}

The weight $T$ now bounds the upper limit of the normalized section
weights in the chosen degrees. We next compare its divisorial coefficient with the minimal
vanishing orders of those complete section spaces. This comparison
holds on every fiber satisfying the stated generic conditions;
the final choice of a fiber on which the integral in
\eqref{ca:global-integral} is finite is not needed for the next lemma.

\begin{lemma}\label{ca:jet}
Let \(D\) satisfy the simultaneous generic conditions above, and fix
a component \(S=D_i\) of \(E\).  Suppose the degrees chosen in
Lemma~\ref{ca:bergman} have minimal actual
vanishing orders
\[
 j_k=\min\{\ord_S(U):0\ne U\in
                  H^0(D,kK_D+r_kH|_D)\}
\]
satisfying \(j_k/k\to0\).  The weight \(T\) constructed in
\eqref{ca:limiting-weight} has zero generic divisorial Lelong
coefficient along \(S\) after restriction to \(D\).
\end{lemma}

If $T|_D$ had positive divisorial coefficient along $S$, an adjoint
estimate would produce a section with the same leading coefficient
and a strictly smaller $L^{2/k}$ integral. We prove both the
preservation of that coefficient and a gain proportional to $1/k$;
the dependence on $k$ matters because the exponent $2/k$ tends to zero.

\begin{proof}
Suppose that coefficient is \(\nu>0\), with Lelong coefficients
normalized by \(\log|z|^2\).  For each \(k\), choose a nonzero
leading section in the image of
\[
 H^0(D,kK_D+r_kH|_D-j_kS)
 \longrightarrow
 H^0\bigl(S,(kK_D+r_kH|_D-j_kS)|_S\bigr).
\]
Among sections with this prescribed image, minimize the
\(2/k\)-integral in \eqref{ca:normalization}.  This minimum exists:
the constraint is a nonempty closed affine subset of a
finite-dimensional vector space, and the integral is continuous,
positive away from zero, and homogeneous of positive degree.
It is therefore coercive.  Denote a minimizer by \(U_k\), and
rescale both it and the prescribed image so that
\eqref{ca:normalization} holds.  All lower coefficients along \(S\)
are zero by the definition of \(j_k\).

Choose \(c>0\) so large that, for all sufficiently large \(k\),
\begin{equation}\label{ca:coefficient-choice}
 c\nu-(1+c)\frac{j_k}{k}>1.
\end{equation}
Choose \(p_0>0\) sufficiently small that
\begin{equation}\label{ca:fiber-integral}
 \int_D |\omega|^2e^{(1+p_0)\psi}<\infty.
\end{equation}
This follows already from the local upper bound for \(T|_D\): the
integrand is bounded by a constant times \(|\omega|^{-2p_0}\),
which is integrable when \(p_0e_j<1\) for every component of the
simple normal crossing divisor \(E\).

Fix a smooth convex function \(g:\R\to\R_{\geq0}\) with
\(0\leq g'\leq1\), equal to \(0\) sufficiently far to the left
and to its argument sufficiently far to the right, and with
\(g''>0\) on \([-1,1]\).  Choose a smooth function
\(0\leq\theta\leq1\) that is zero on \((-\infty,-1]\), one
on \([1,\infty)\), and whose derivative is supported in
\([-1,1]\).  For \(b>0\), put
\begin{equation}\label{ca:adjoint-weight}
 \Phi=(k-1-c-p_0)\varphi_k+cT|_D
          +p_0\bigl(\varphi_0+g(x_k-b)\bigr)+r_k\varphi_H.
\end{equation}
This is a weight on the integral bundle
\((k-1)K_D+r_kH|_D\): the coefficients of its canonical weights
sum to \(k-1\), and \(x_k\) is a global function off the divisors.
We henceforth take \(k>1+c+p_0\).  The constants $c,p_0$ are fixed.
For each fixed $k$ and cutoff $b$, we first remove the regularization
used below.  We then estimate the upper limit as $k\to\infty$ with
$b$ fixed, and finally let $b\to\infty$.  The contradiction will use
one sufficiently large $b$ followed by sufficiently large $k$.

\smallskip\noindent\textbf{The curvature and the adjoint estimate.}\enspace
Off \(E\) and the zero divisor of \(U_k\), one has
\(\ddc\varphi_k=-(r_k/k)h\), \(\ddc\varphi_0=0\), and
\(\ddc x_k=-(r_k/k)h\).  Direct differentiation of
\eqref{ca:adjoint-weight} gives
\begin{align}
 \ddc\Phi={}&\frac{r_k}{k}(1+c+p_0-p_0g'(x_k-b))h
       +c\ddc(T|_D) \notag\\
 &+p_0g''(x_k-b)\frac{i}{2\pi}
               \partial x_k\wedge\bar\partial x_k.
                                                   \label{ca:curvature}
\end{align}
In particular it dominates \((1+c)r_kh/k\) and retains the displayed
gradient term.  The resulting adjoint estimate produces
\(V_k\in H^0(D,kK_D+r_kH|_D)\) such that, on the complement of
the divisors,
\[
 u_k=\theta(x_k-b)U_k-V_k
\]
satisfies
\begin{equation}\label{ca:l2-estimate}
 \int_D\left|\frac{u_k}{U_k}\right|^2w_k
 \leq C_{p_0}\int_{\{|x_k-b|\leq1\}}w_k,
 \qquad
 w_k=|\omega|^2
 e^{(1+c+p_0)x_k-c\psi-p_0g(x_k-b)}.
\end{equation}
The constant is independent of \(b\) and \(k\): the retained gradient
term controls the datum on the cutoff shell, even as the ample
coefficient tends to zero. We give the
singular-weight justification, including the choice of complete
metric, before using this inequality.

Subtract a fixed smooth reference metric from \(T|_D\) and apply
Theorem~14.12(a),(c) and Corollary~14.13(c) of \cite{Demailly12}.  We obtain
weights \(T_\ell\) converging almost everywhere to \(T|_D\),
uniformly bounded above locally, smooth outside an analytic set
\(Z_\ell\), and satisfying
\(\ddc T_\ell\geq-\epsilon_\ell h\), where
\(\epsilon_\ell\to0\).  The cited construction gives pointwise
approximation above the original potential.  In any event, taking
a smooth regularized maximum with \(\varphi_0\), of width tending
to zero, ensures \(T_\ell\geq\varphi_0\).  The regularized
maximum is invariant under simultaneous translation of its
arguments, so it gives a weight on the same bundle.  Its
monotonicity and convexity preserve the lower curvature bound.
It is smooth off \(Z_\ell\cup E\); preservation of analytic
singularities of the maximum itself is not required.

Keep \(k,b\) fixed and replace \(T|_D\) by \(T_\ell\) in
\eqref{ca:adjoint-weight}, obtaining \(\Phi_\ell\).
For large \(\ell\), the loss \(c\epsilon_\ell h\) is less than
half the ample term in \eqref{ca:curvature}.  On
\[
 D\setminus\bigl(Z_\ell\cup E\cup\operatorname{div}(U_k)\bigr)
\]
the weight is smooth with positive curvature and the same gradient
lower bound.  This complement admits a complete K\"ahler metric.
Indeed its deleted analytic set is the zero set of a section of a
Hermitian vector bundle: twist its coherent ideal by a sufficiently
ample bundle and use finitely many generators.  Lemma~12.9 of
\cite{Demailly12} applies, because \(D\) is compact K\"ahler and
the curvature of that smooth bundle has a uniform upper bound.

Write \(A_\ell=i\partial\bar\partial\Phi_\ell\) for the full
positive curvature form on the complement.  Add \(A_\ell\) to a
complete K\"ahler metric there, obtaining a complete metric
\(\widehat h_\ell\geq A_\ell\).  This particular choice also
verifies the ordinary \(L^2\) hypothesis for the datum
\[
 f_k=\theta'(x_k-b)\bar\partial x_k\,U_k.
\]
To see both this and the uniform estimate explicitly, let
\(G\) be the matrix of an auxiliary K\"ahler metric and \(A\)
the positive curvature matrix.  For a bundle-valued \((n,1)\)-form
with local coefficient \(U\eta\), matrix cancellation gives
\begin{equation}\label{ca:metric-cancellation}
 \bigl\langle[A,\Lambda_G]^{-1}(U\eta),U\eta\bigr\rangle_G
       e^{-\Phi}\,dV_G
 = |U|^2e^{-\Phi}\,\eta^*A^{-1}\eta\,d\lambda,
\end{equation}
up to the fixed convention for the canonical volume form.
The determinants from the top holomorphic degree and the volume
cancel, and so do the metric factors in the antiholomorphic degree
and in the curvature operator.  The solution norm in bidegree
\((n,0)\) is likewise independent of \(G\).  When \(G\geq A\),
the ordinary datum norm is bounded by the right side of
\eqref{ca:metric-cancellation}.

The gradient term in \eqref{ca:curvature} implies, for
\(\eta=\bar\partial x_k\), the matrix bound
\[
 \eta^*A_\ell^{-1}\eta
       \leq\frac{1}{p_0g''(x_k-b)}.
\]
Since \(\theta'\) is supported where \(g''\) has a positive
minimum, the inverse-curvature density of \(f_k\) is at most a
constant times the shell density.  That density is integrable.
The datum is \(\bar\partial\)-closed, being the derivative of
\(\theta(x_k-b)U_k\) on the complement.  The complete K\"ahler
estimate, Theorem~5.1 of \cite{Demailly12}, consequently gives
\(\bar\partial u_{k,\ell}=f_k\) with the desired bound for
\(\Phi_\ell\).

Here the local adjoint density has precisely the required
exponents:
\begin{align*}
 |U_k|^2e^{-\Phi_\ell}
 &=\exp\bigl((1+c+p_0)\varphi_k-cT_\ell
                 -p_0(\varphi_0+g(x_k-b))\bigr)\,d\lambda\\
 &=|\omega|^2
       e^{(1+c+p_0)x_k-c(T_\ell-\varphi_0)-p_0g(x_k-b)}.
\end{align*}
On the fixed shell, it is bounded above by
\(e^{(1+c+p_0)(b+1)}|\omega|^2\), because
\(T_\ell-\varphi_0\geq0\).  Moreover \(\Phi_\ell\) is bounded
above locally uniformly in \(\ell\), using the positive coefficient
\(k-1-c-p_0\) and
\[
 \varphi_0+g(\varphi_k-\varphi_0-b)
       \leq\max\{\varphi_0,\varphi_k-b\}+C_g.
\]
The weighted solution estimate therefore bounds the ordinary local
\(L^2\) norms of \(u_{k,\ell}\).  The holomorphic section
\(V_{k,\ell}=\theta(x_k-b)U_k-u_{k,\ell}\) extends across its
deleted analytic set by the distributional \(L^2\) removable-singularity
lemma \cite[Lemma~12.10]{Demailly12}. In a local bundle frame the
right side of the distributional $\bar\partial$ equation is zero,
so the extended coefficients have holomorphic representatives.
The extensions lie in the fixed finite-dimensional space
\(H^0(D,kK_D+r_kH|_D)\), with bounded ordinary norm.  Take a
subsequence converging as global sections.  Fatou's lemma for the
left side and dominated convergence on the fixed shell give
\eqref{ca:l2-estimate}.  For this passage one may discard the
countable union of the sets \(Z_\ell\); no nesting of those sets is
needed.

\smallskip\noindent\textbf{The shell becomes small.}\enspace
Put \(B=1+c+p_0\).  For fixed \(b\), the shell densities in
\eqref{ca:l2-estimate} are bounded by
\(e^{B(b+1)}|\omega|^2\).  At a point off \(E\), a subsequence
contributing to their limsup has a further subsequence with
\(x_k\to x_*\in[b-1,b+1]\).  Lemma~\ref{ca:bergman} gives
\(x_*\leq\psi\); thus \(\psi\geq b-1\) and
\[
 Bx_*-c\psi-p_0g(x_*-b)\leq(1+p_0)\psi.
\]
Reverse Fatou therefore yields
\begin{equation}\label{ca:shell-tail}
 \limsup_{k\to\infty}\int_{\{|x_k-b|\leq1\}}w_k
 \leq\int_{\{\psi\geq b-1\}}|\omega|^2e^{(1+p_0)\psi}.
\end{equation}
By \eqref{ca:fiber-integral}, the right side tends to zero as
\(b\to\infty\).

\smallskip\noindent\textbf{The prescribed leading section is preserved.}\enspace
Near a general point of \(S=(z=0)\), write
\(U_k=z^{j_k}a_k\) and \(\omega=z^{e_i}a_0\), with the
coefficients nonvanishing in the chosen frames.  Since
\(j_k/k<e_i\), the function \(x_k\) tends to plus infinity as
\(z\to0\).  Thus \(\theta(x_k-b)=1\) and
\(g(x_k-b)=x_k-b\) sufficiently near that point.
The divisorial decomposition of the positive curvature current of
\(T|_D\) gives
\(T|_D\leq\nu\log|z|^2+O(1)\): subtracting
\(\nu[S]\) leaves a positive current, whose local potential is
bounded above.  See \cite[Chapter~III, Proposition~8.16]{DemaillyCADG}.
Consequently
\begin{equation}\label{ca:jet-weight-order}
 \Phi\leq\left((k-1-c)\frac{j_k}{k}+c\nu\right)
                      \log|z|^2+O(1).
\end{equation}
The coefficient exceeds \(j_k+1\) by
\eqref{ca:coefficient-choice}.  In this neighborhood
\(u_k=U_k-V_k\) is holomorphic.  If its vanishing order were at
most \(j_k\), the transverse integral against \(e^{-\Phi}\)
would diverge, contradicting \eqref{ca:l2-estimate}.  Hence
\(U_k-V_k\) vanishes along \(S\) to order at least \(j_k+1\).
This divisorial statement says exactly that \(V_k\) has the
prescribed leading section and all the prescribed zero lower
coefficients.

\smallskip\noindent\textbf{A gain of order \(1/k\).}\enspace
Let \(\rho_k\,dV\) be the normalized density in
\eqref{ca:normalization} relative to a fixed smooth volume on
\(D\), and put \(q_k=|u_k/U_k|\) off the zero divisor.
The functions \(\rho_k\) have a uniform upper bound by
Lemma~\ref{ca:bergman}. In local canonical frames the density is
\[
 w_k=\exp\bigl(B\varphi_k-cT|_D
             -p_0(\varphi_0+g(x_k-b))\bigr)\,d\lambda.
\]
The local upper bounds for \(T|_D\), and the earlier bound
$\varphi_0+g(x_k-b)\leq\max\{\varphi_0,\varphi_k-b\}+C_g$,
therefore give, on a finite cover by smooth frames,
\begin{equation}\label{ca:density-lower-bound}
 w_k\geq C^{-1}\rho_k^B\,dV
\end{equation}
with \(C\) independent of \(k\) and of \(b\) bounded below.
Fix \(0<\eta<\min\{1,2/B\}\).  H\"older's inequality gives
\begin{align}
 \int_D q_k^\eta\rho_k\,dV
 &\leq
 \left(\int_D q_k^2\rho_k^B\,dV\right)^{\eta/2}
 \left(\int_D
       \rho_k^{(2-B\eta)/(2-\eta)}\,dV\right)^{1-\eta/2}
                                                        \label{ca:holder}\\
 &\leq C_\eta
       \left(\int_D q_k^2w_k\right)^{\eta/2}.\notag
\end{align}
The exponent in the second factor is positive, so its uniform
bound follows from compactness and the upper bound for \(\rho_k\).
Equations~\eqref{ca:l2-estimate} and \eqref{ca:shell-tail} show
that this moment can be made arbitrarily small by first fixing
\(b\) sufficiently large and then taking \(k\) sufficiently large.

The set where \(\theta(x_k-b)\ne0\) is contained in
\(\{\varphi_0\leq C-b+1\}\), by the upper bound for
\(\varphi_k\).  Its smooth volume tends to zero with \(b\),
and therefore its \(\rho_k\)-mass tends to zero uniformly in
\(k\).  Markov's inequality applied to
\eqref{ca:holder} now gives a set of \(\rho_k\)-mass at least
\(1/2\) on which \(\theta(x_k-b)=0\) and \(q_k\leq1/2\).
On this set \(|V_k/U_k|=q_k\leq1/2\); everywhere else,
\(|V_k/U_k|\leq1+q_k\).

For \(0<p\leq\eta/2\) and \(q\geq0\),
\begin{equation}\label{ca:small-exponent}
 (1+q)^p\leq1+C_\eta p q^\eta.
\end{equation}
For \(q\leq1\), this follows from concavity.  For \(q\geq1\),
integrate the derivative with respect to the exponent and use
\((1+q)^{\eta/2}\log(1+q)\leq C_\eta q^\eta\).
Taking \(p=2/k\), we therefore obtain
\begin{align}
 \int_D |V_k|^{2/k}e^{-r_k\varphi_H/k}
 &=\int_D |V_k/U_k|^p\rho_k\,dV\notag\\
 &\leq1-\frac12(1-2^{-p})
             +C_\eta p\int_D q_k^\eta\rho_k\,dV<1
                                                        \label{ca:strict-gain}
\end{align}
for the indicated choices of \(b\) and then \(k\).
Indeed \((1-2^{-p})/p\to\log2\), whereas the moment can be
made arbitrarily small independently of large \(k\).
This is a strict improvement of the minimizing integral for a
section with the same prescription, a contradiction.  Therefore
\(\nu=0\).
\end{proof}

\subsection{Excluding poles with zero limiting fixed order}

We now apply the estimates to a horizontal pole of $v$.
The degrees and the limiting weight will be chosen from the
geometric generic section spaces. Only after applying
Lemma~\ref{ca:integrability} to that fixed weight will we choose
a complex fiber and the minimizing sections on it.

\begin{proposition}\label{ca:poles}
Every horizontal pole divisor of \(v\) has positive limiting fixed
order on a simultaneous very general fiber.  In particular, for
one sufficiently small positive rational \(\epsilon\), every such
divisor has positive actual asymptotic fixed order for
\(K_D+\epsilon H|_D\).
\end{proposition}

\begin{proof}
Suppose a horizontal pole divisor induces \(S=D_i\) and has
zero limiting fixed order on the geometric generic fiber.  Choose
rational perturbations tending to zero and degrees in which the
normalized minimal orders approach their infima.  Passing to powers
as necessary gives integers
\(k\to\infty\), \(r_k>0\), with
\[
 \frac{r_k}{k}\longrightarrow0,
 \qquad \frac{j_k}{k}\longrightarrow0,
\]
where \(j_k\) is the minimal actual order of the whole section
space.  Use the corresponding relative section spaces and vanishing
filtrations on the fixed prepared family, and impose the countably many
base-change conditions that preserve these orders.  No further field
extension or spreading of individual sections is required.
Lemma~\ref{ca:bergman} now constructs one global weight $T$ from the
chosen degrees, before any final fiber is selected.

Apply Lemma~\ref{ca:integrability} to this $T$.  Fubini gives a
full-measure set of good fibers on which the integral is finite.
The excluded algebraic conditions form a countable subset of the
complex curve and hence have measure zero.  Choose $D$ satisfying
both requirements.  Lemma~\ref{ca:jet} applies on this same fiber
and proves that $T|_D$ has zero generic divisorial coefficient at $S$.
The normalized minimizing sections in that Lemma are chosen only now;
they do not change $T$.

In a smooth chart over the base, contraction gives
\[
 \Theta=\iota_v(dt\wedge\omega).
\]
The density of \(i^{n^2}h\wedge\Theta\wedge\overline\Theta\)
relative to base area is comparable, by the smooth positive metric,
to \(|v|_h^2|\omega|^2\).  This formula includes the mixed
components of the metric.  At a generic point of a pole of order
at least one, shrink the chart so its leading coefficient is
bounded below.  Since \(\ord_S(\omega)=e_i\), the finite fiber
integral from \eqref{ca:global-integral} would then dominate a
positive constant times
\begin{equation}\label{ca:pole-integral}
 \int |z|^{-2-2\delta e_i}e^{(1+\delta)T|_D}\,d\lambda,
 \qquad S=(z=0).
\end{equation}
We show that this last integral is infinite.

The generic Lelong coefficient is zero.  The analyticity of Lelong
upper-level sets, or equivalently the generic-Lelong-number
statement in \cite[Chapter~III, Lemma~8.15]{DemaillyCADG}, implies
that the point Lelong number is zero at almost every center on a
small smooth patch of \(S\).  For a local plurisubharmonic weight
\(T\), the logarithmic mean characterization
\cite[Chapter~III, Corollary~6.8 and Example~6.9]{DemaillyCADG}
then gives
\[
 \frac{1}{\vol(B(a,r))}\int_{B(a,r)}T\,d\lambda
                   =o(|\log r|)
\]
at almost every such center \(a\).  The characterization for ball
means follows from that for sphere means by radial integration.

Fix \(\varepsilon>0\).  For almost every center, the ball mean
is at least \(\varepsilon\log r\) at all sufficiently small
dyadic radii.  The sets on which this holds beyond one fixed
dyadic index have a union of full divisor measure.
Hence there is a measurable set \(A\) of positive divisor measure
on a compact patch of \(S\), and one common starting index, for
which the lower bound holds uniformly.  In dimension one the
divisor measure here is counting measure.

On \(B(a,r)\), \(|z|\leq Cr\).  Jensen's inequality therefore
gives, for the nonnegative integrand \(F\) in
\eqref{ca:pole-integral},
\[
 \int_{B(a,r)}F\,d\lambda
 \geq C^{-1}r^{2n-2-2\delta e_i+(1+\delta)\varepsilon}
 \qquad(a\in A).
\]
Integrate this inequality in \(a\).  For each ambient point, the
measure of centers \(a\in A\) whose radius-\(r\) ball contains
it is at most \(Cr^{2n-2}\), by projection to the smooth
tangential coordinates of \(S\).  Fubini consequently yields
\[
 \int F\,d\lambda
 \geq C^{-1}r^{-2\delta e_i+(1+\delta)\varepsilon}.
\]
The overlap bound is constant when \(n=1\), so the same formula
holds in that case.  Since \(e_i\geq1\), choose
\(\varepsilon<2\delta e_i/(1+\delta)\) and let the dyadic
radius tend to zero.  The right side diverges, contradicting
\eqref{ca:pole-integral} and the Fubini consequence of
\eqref{ca:global-integral}.  A pole with zero limiting fixed order
is therefore impossible.

There are only finitely many horizontal pole divisors.  For each,
positivity of its limit and monotonicity give a sufficiently small
rational perturbation with positive actual asymptotic fixed order.
Taking one common smaller rational perturbation proves the last
assertion.
\end{proof}


\subsection{Contraction of poles and the geometric generic fiber}

\begin{proof}[Proof of Theorem~\ref{cc:criterion}]
We may assume $n>0$, and use the preparation and notation of
Lemmas~\ref{cc:preparation}--\ref{cc:zeros}.
Proposition~\ref{ca:poles} gives one sufficiently small positive rational
$\varepsilon$ for which every horizontal pole divisor has positive
actual asymptotic fixed order for $K_D+\varepsilon H|_D$.
We fix this same perturbation for the relative model construction.
The order comparisons hold on the geometric generic fiber and on
simultaneous very general fibers by base change in the countably many
divisible section degrees.

Over the generic curve choose an effective rational divisor
$\Delta\sim_{\Q}\varepsilon H$ such that the pair is klt.  One obtains
it by taking a general sufficiently divisible member of the ample
system and a small coefficient.  In the equivariant case average its
translates; the klt condition is preserved because discrepancies are
linear in the boundary and the translates are klt.  Spread this choice
and shrink the curve so that the total pair is klt.  The relative
adjoint divisor is big: on a good fiber $K_D\sim E\ge0$, and an ample
positive perturbation of an effective divisor is big.

Apply \cite[Theorem~1.2(2)]{BCHM} to this projective morphism over the
shrunken curve.  Its hypotheses are a klt pair and a big relative
adjoint divisor.  It gives a normal relative log canonical model
\[
                         \chi:P\dashrightarrow Q
\]
with $Q$ projective over the curve.  Since the adjoint divisor is
relatively big, this map is birational.  The ample-model comparison
\cite[Definitions~3.6.4 and~3.6.6, Lemma~3.6.5(4)]{BCHM} shows that it
is a birational contraction, so it extracts no divisors.  In a common
resolution the pullback adjoint system is the pullback of the ample
system on $Q$ plus an effective exceptional fixed part.

Every horizontal pole divisor is contracted by $\chi$.  Indeed a prime
which survived would be tested at its corresponding prime on the ample
model.  Sufficiently divisible ample systems have no fixed component
there, and the exceptional section comparison changes no order at that
prime.  Its asymptotic fixed order would therefore be zero, contrary to
the choice of $\varepsilon$.
The canonical ring is invariant under the finite group when $\Delta$
is invariant, and hence the action descends to $Q$.

Choose a nonvanishing algebraic vector field on a smaller open of the
curve.  Normalize the kernel line of $\Theta$ against that vector
field, and denote the resulting rational derivation of $\C(P)=\C(Q)$
by $\partial$.  It is rational because $\Theta$ is algebraic and its
kernel and prescribed base projection are given by algebraic linear
equations.  This normalization multiplies the previous local $v$ by a
nonzero base function, so it creates no horizontal pole.  Remove the
finitely many vertical pole loci and the zeros or poles of the chosen
base vector field.

The derivation is regular at every prime divisor of $Q$.  Because
$\chi$ extracts no divisors, each such prime corresponds to a prime on
$P$, with the same discrete valuation ring; it cannot be one of the
contracted pole divisors.  It remains to pass through codimension two.
For any normal affine open with coordinate ring $A$, and any $f\in A$,
regularity at the height-one points gives
\[
           \partial f\in\bigcap_{\operatorname{ht}\mathfrak p=1}
                                      A_{\mathfrak p}=A.
\]
Thus $\partial$ is a regular derivation on the entire normal space.
Its projection to the base is still the chosen base vector field,
since this equality is an identity of derivations on the function
field.  Locally reparametrize the base analytically so that
$\partial t=1$.

A regular derivation on a complex analytic space has local holomorphic
flows, including at singular points.  Here is the local justification.
Embed the space in a polydisk with coordinate functions $z_1,\ldots,z_N$,
lift the functions $\partial z_j$ to holomorphic functions on the
polydisk, and form the corresponding ambient vector field.  It
preserves the defining ideal, because its induced derivation on the
quotient is $\partial$.  Its flow preserves the analytic subspace:
for finitely many generators of the ideal, their values along an
integral curve satisfy a homogeneous linear differential system with
zero initial value.  Uniqueness makes the restrictions of these local
flows independent of the ambient lifts and permits them to glue.

Properness now gives a uniform flow time near a compact fiber.  More
explicitly, the inverse image of a sufficiently small closed base disk
is compact; a finite cover by flow neighborhoods provides a common
smaller disk of complex times.  Since $\partial t=1$, forward and backward
flows identify the fiber over its center with every fiber over a
smaller disk.  These identifications are biholomorphic, and hence
algebraic because the fibers are projective.  They commute with
$\Gamma$: the kernel of $\Theta$ is invariant, and its normalization by
a base vector field is unique.  Shrink the algebraic open as needed so
that the original family and the model have birational fibers and the
model fibers under consideration are geometrically integral.

Fix one fiber $Q_{t_0}$ in this disk.  Graphs of isomorphisms from
$Q_{t_0}$ to fibers of $Q$ belong to countably many finite-type loci of
relative Hilbert schemes.  Requiring equivariance is a closed condition
on these graphs.  Their images in the algebraic curve are constructible
and together contain the disk, by the algebraic flow identifications.
If none were dominant, each image would be a finite subset of the
curve; their countable union could not contain a disk.  Some graph
locus is therefore dominant.  Its generic fiber has a closed point
defined over a finite extension of the curve's function field.  Over
that extension the generic fiber of $Q$ is isomorphic to the fixed
$Q_{t_0}$, equivariantly when prescribed.

Take a smooth projective resolution $D_0$ of $Q_{t_0}$, equivariant in
the finite-group case.  The birational comparison with $P$ gives the
required birational identification of its generic fiber with $D_0$
after a finite extension.  The preliminary finite changes of the curve
compose with this extension.  For $n=0$ the fiber is a point and the
conclusion is immediate.  This completes the proof.
\end{proof}

\subsection{From curves to geometric birational constancy}

The curve criterion gives birational identifications after finite
extensions of curve function fields. We now use algebraic loci of
graphs to obtain the corresponding statement on an arbitrary
constant-line locus and on a geometric generic parameter fiber.

\begin{corollary}[Constant-line loci]
\label{cc:loci}
\label{cor:period-loci}
Let $\pi:\mathcal P\to B$ be a smooth projective morphism with
connected fibers over a smooth connected complex algebraic variety.
Let $T\subseteq B$ be a connected smooth algebraic locus on which the
very general fiber $D$ satisfies $\kappa(D)=0$ and $h^0(D,K_D)=1$.
If the entire ordinary top Hodge line of the restricted family has
constant projective period in a flat reference, its geometric generic
fiber is birational to a smooth projective variety over $\C$.
The conclusion is equivariant for any prescribed finite fiberwise
birational group action.

More generally, let $B\to Z$ be a morphism of smooth connected complex
algebraic varieties with geometrically integral generic fiber.  If the
same fiber and constant-line hypotheses hold on a nonempty smooth
open of every very general fiber over $Z$, then the family over the
geometric generic fiber of $B\to Z$ is birationally constant over
$\overline{\C(Z)}$, with the prescribed finite action.  These
birational identifications remain valid after algebraically closed
extension of their fields of definition.
\end{corollary}

\begin{proof}
Restrict the integral polarized cohomology variation to $T$.  Its
entire highest Hodge step has rank one.  Constancy of the projective
line makes it a flat rank-one local subsystem, and
Corollary~\ref{hodge:restricted-character}, applied anew to this
restricted variation, gives finite character.  On an algebraic curve
in $T$, Theorem~\ref{cc:criterion} therefore applies, including its
equivariant assertion.

We pass from curves to the geometric generic point of $T$.  After
shrinking $T$ to an affine open, graphs of birational identifications between pairs of
fibers are parametrized by countably many finite-type algebraic loci
over $T\times T$.  For example, take relative Hilbert schemes of the
product family and their constructible loci of geometrically integral
graphs on which both projections are dominant of generic degree one.
Geometric integrality and these degree conditions are imposed after
flat stratification.  Equivariance is
expressed by invariance of the graph under the diagonal finite action;
a birational action is first regularized as in
Lemma~\ref{cc:preparation}.  The images of the graph loci in
$T\times T$ are constructible.

For $\dim T=0$ the conclusion is immediate, and for $\dim T=1$ it is
Theorem~\ref{cc:criterion}.  Suppose $\dim T\ge2$.  Choose a smooth
projective compactification and a sufficiently ample complete-intersection
family of curves.  The incidence of a curve with two distinct points
on it has an irreducible open dominating $T\times T$: sufficiently
ample systems separate the two point conditions, and their general
members through those points are smooth and connected.  Its general
fiber over the curve parameter space is the irreducible surface of
pairs of points on that curve.

If none of the graph images were dense in $T\times T$, their closed
image closures would pull back to countably many proper closed subsets
of this incidence.  For a very general curve, each such subset meets
its surface of pairs in a proper closed subset; those not dominating
the curve parameter space are avoided.  The curve can also satisfy
the countably many generic section conditions giving $\kappa(D)=0$
and $h^0(D,K_D)=1$.  Theorem~\ref{cc:criterion} and spreading its
birational identification give a dense open of pairs of fibers on
that curve which are birational, equivariantly when required.  A
countable union of proper closed subsets of an irreducible complex
surface cannot contain this open.  This contradiction proves that
some graph locus dominates $T\times T$.

A dense constructible image contains a nonempty open.  Specialize the
first coordinate to a sufficiently general complex point so that the
slice of this open is dense in the second factor.  The first fiber is
now a fixed smooth projective variety.  The graph parameter space over
this slice dominates $T$, and a closed point of its generic fiber is
defined over a finite extension of $\C(T)$.  It gives the desired
birational identification, with the finite action.

For the relative assertion use these same graph loci over
$B\times_Z B$, taking the component that dominates $Z$.  Its geometric
generic fiber is integral by hypothesis.  If no graph locus dominated
this component, the closures of their images would remain proper on
every very general fiber over $Z$, after excluding countably many
proper specialization loci.  This contradicts the assertion just
proved on those complex fibers.  Thus one graph locus dominates over
the generic point of $Z$.  Over $b_Z=\overline{\C(Z)}$, its
constructible image contains an open in the space of pairs.  Choose a
$b_Z$-point in a suitable open for the first projection.  Its fiber is
a fixed smooth projective $b_Z$-variety, and the graph locus on the
remaining slice supplies the birational identification after a finite
extension of the function field of the geometric generic $B$-fiber.
This proves the relative conclusion.  Finally, all the graphs, inverse
maps, and equivariance identities are algebraic data, so these
identifications persist under field extension.
\end{proof}

\begin{proposition}[Constancy in the $h$-fiber directions]
\label{ro:constancy}
For very general $z\in Z$, the geometric generic fiber of the
$J$-family over $W_z$ is birational to the base change of a smooth
projective variety over $\C$.  More intrinsically, after extending the
function field of $Z$ to its algebraic closure, the geometric generic
fiber of $x$ over the resulting $h$-fiber is birationally constant over
that algebraically closed field.
\end{proposition}

\begin{proof}
On the smooth family locus $W^\circ$, Lemma~\ref{ro:root} gives a
smooth projective cyclic-cover family with connected fibers,
Kodaira dimension zero, and a one-dimensional \emph{entire}
ordinary top-form space. On the connected open of every very general
$W_z$, Corollary~\ref{ro:flat-character} makes its projective top line
constant. Corollary~\ref{cc:loci} therefore makes that restricted
cover family geometrically birationally constant, equivariantly for
its cyclic deck group. Taking invariant function fields gives the
asserted constancy of the original $J$-family.

For the relative assertion, apply the second part of
Corollary~\ref{cc:loci} to $W^\circ\to Z$, shrinking $Z$ if necessary.
Its geometric generic fiber is integral by
Proposition~\ref{it:diagram}, and the preceding paragraph verifies
the fiber and constant-line hypotheses on its very general complex
fibers. Over $b_Z=\overline{\C(Z)}$ the reference cover and its finite
deck action are consequently defined over $b_Z$. Their invariant
function field has a smooth projective $b_Z$-model, whose base
extension is birational to the original geometric generic $x$-fiber.
This is geometric birational constancy; the descent of a chosen
trivialization and of the whole original fiber is addressed in the
remaining sections.
\end{proof}
